\documentclass[10pt,a4paper]{amsart}
\usepackage[margin=19mm]{geometry}
\usepackage{amsmath,amssymb,mathtools}
\usepackage[scr=rsfs,cal=euler]{mathalfa}
\usepackage{xfrac}
\usepackage[unicode,hidelinks,bookmarks=false]{hyperref}
\newcommand{\ie}{i.e.\ }
\newcommand{\cf}{cf.\ }
\newcommand{\resp}{respectively\ }
\newcommand{\Subsection}{\subsection}
\providecommand{\nobreakdash}{\nobreak\hbox{-}}
\setlength{\emergencystretch}{2em}
\multlinegap0pt
\usepackage{graphics}
\usepackage{epsfig}
\usepackage{colortbl}
\newcommand{\bZ}{\mathbb{Z}}
\newcommand{\bR}{\mathbb{R}}
\newcommand{\bC}{\mathbb{C}}
\newcommand{\e}{{\rm e}}
\newcommand{\dd}{\, {\rm d}}
\setlength{\topmargin}{-1.5cm} \setlength{\oddsidemargin}{0.2cm}
\setlength{\evensidemargin}{-0.2cm}
\setlength{\textheight}{22.5cm} \setlength{\textwidth}{15.5cm}
\setlength{\headsep}{1.4cm}
\newtheorem{theorem}{Theorem}
\newtheorem{proposition}[theorem]{Proposition}
\newtheorem{lemma}[theorem]{Lemma}
\newtheorem{remark}[theorem]{Remark}
\newtheorem{corollary}[theorem]{Corollary}
\def\qed{\hfill\begin{picture}(6.5,6.5)%
\put(0,0){\framebox(6.5,6.5){}}\end{picture}}
\begin{document}
\title[Splitting methods for nonlinear Schr"odinger equation withou]{Splitting methods for nonlinear Schr\"odinger equation without order reduction}
\author{Facultad de Ciencias, Universidad de Valladolid,; Paseo de Belén 7, 47011 Valladolid,; Spain ạte; Carlos Arranz-Simón; Begoña Cano}
\date{}
\maketitle
\noindent\emph{Source and licence.} Splitting methods for nonlinear Schr\"odinger equation without order reduction. Version arXiv:2607.08387v1 (CC BY 4.0). This material is distributed under the Creative Commons Attribution 4.0 International licence, http://creativecommons.org/licenses/by/4.0/, as recorded in the arXiv OAI licence metadata for this exact version and shown on the abstract page https://arxiv.org/abs/2607.08387v1. Full rights notice: licenses/arxiv-2607.08387-CC-BY-4.0.txt.\par\smallskip
\noindent\textit{Editorial scope.} Counted unit: the original \texttt{theorem} environment \texttt{reg2} at source lines 253--269 together with its complete proof at lines 271--386; the delivered document keeps source lines 94--387 so that every referenced label and every citation resolves inside it. Native editable source is the paper's own concatenated TeX; the AMS wrapper is editorial formatting only. No publisher class, upstream script or unrelated artwork is redistributed.
\begin{eqnarray}
\dot{u}(t)&=& i \big(A u(t)+f(|u(t)|^2) u\big), \quad 0 \le t \le T,\nonumber \\
u(0)&=&u_0,  \nonumber \\
\partial u(t)&=& g(t), \label{nlse}
\end{eqnarray}
where $A: H^2(\Omega)\to L^2(\Omega)$ denotes the operator which applies the second derivative in space in one dimension and the Laplacian in general for a bounded domain $\Omega$ with regular enough boundary. On the other hand, $f$ is a smooth enough real function, $u_0\in H^2(\Omega)$ and $\partial: H^2(\Omega) \to H^ {\frac{3}{2}}(\partial \Omega)$ is the Dirichlet operator.

In the following, we will use the operators $A_0=A|_{\mbox{ker}(\partial)}$ and $K: H^{\frac{3}{2}}(\partial \Omega) \to H^2(\Omega)$, so that $K g$ is the solution of the elliptic problem
\begin{eqnarray}
A u &=&0, \nonumber \\
\partial u&=&g, \nonumber
\end{eqnarray}
which is well-known to satisfy \cite{ACP}
\begin{eqnarray}
\|Kg \|_{L^2(\Omega)}\le C \|g\|_{H^{\frac{3}{2}}(\partial \Omega)},
\nonumber
\end{eqnarray}
for some constant $C$.
Moreover, if $g \in H^{\frac{3}{2}}(\partial \Omega)\cap L^{\infty}(\Omega)$, for some other constant it happens that \cite{S}
\begin{eqnarray}
\|Kg \|_{L^\infty(\Omega)}\le C \|g\|_{L^\infty(\partial \Omega)}.
\label{acotK2}
\end{eqnarray}
It is well-known that advancing a stepsize $\tau$ with Strang method to solve a general differential system of the form
$$\dot{u}=f_1(u)+f_2(u),$$
consists of the following
$$u_{n+1}=\phi_{\frac{\tau}{2}}^{f_2}\circ\phi_{\tau}^{f_1}\circ\phi_{\frac{\tau}{2}}^{f_2}(u^n),$$
where $\phi_\tau^{f_1}$ and $\phi_\tau^{f_2}$ are the exact flows after time $\tau$ of the respective systems
$$
\dot{u}=f_1(u), \quad \dot{u}=f_2(u).$$
In the case of
(\ref{nlse}), one of the terms in the splitting would be
\begin{eqnarray}
\dot{u}&=& i f(|u|^2) u, \label{nolinear}
\end{eqnarray}
which can be exactly solved because $|u|^2$ remains constant under this evolution \cite{WH}
$$
\frac{d}{dt} |u|^2= \dot{u} \bar{u}+ u \dot{\bar{u}}=i f(|u|^2) u  \bar{u}-i u f(|u|^2) \bar{u}=0,$$
so this can be solved as a linear system. More precisely, the solution of (\ref{nolinear}) from an initial condition $u_0$ would be
$$
u(t)=e^{i t f(|u_0|^2)}u_0.$$

On the other hand, for the solution of the other term of the splitting
\begin{eqnarray}
\dot{u}&=& i A u, \label{linear}
\end{eqnarray}
to be completely determined, apart from an initial condition, a boundary condition is also necessary. In that sense, in \cite{CR}, some boundaries are suggested which are of course related to the boundary $g$ in (\ref{nlse}) and for which no order reduction is observed. More precisely, second global order is observed when applying Strang method with those boundaries and solving exactly its various stages. If an analytic expression for $g$ is known, those boundaries can be calculated without resorting to numerical differentiation, although just local order $2$ is achieved and a summation-by-parts argument must be applied to justify the same global order. In any case, $\dot{g}$ is necessary and, if no analytic expression of $g$ is available, numerical differentiation would also be required. Finally, notice that, as justified in Remark 1 of \cite{CR}, if local order $3$ is searched for, numerical differentiation in space from the numerical solution and not directly from the data in (\ref{nlse}) would also be required.

\section{Midpoint rule to solve the stiff part of Strang method without order reduction and conserving symmetry}

In this section, without knowing the exact boundary in (\ref{linear}) which would lead to local order $3$ or more, we suggest a way to integrate it (again with local order $3$) just by knowing the exact values of the boundary at the beginning and end of integration of such a stage.
We notice that, at each step, such part of Strang method corresponds to solve
\begin{eqnarray}
\dot{w}^n(s)&=&i A w^n(s), \nonumber \\
\partial w^n(s)&=& g^n(s), \nonumber \\
w^n(0)&=& e^{i \frac{\tau}{2}f(|u_n|^2)} u_n, \label{midstrang}
\end{eqnarray}
where $g^n$ should be a function which satisfies
\begin{eqnarray}
g^n(0)&=& e^{i \frac{\tau}{2}f(|g(t_n)|^2)} g(t_n), \nonumber \\
g^n(\tau)&=& e^{-i \frac{\tau}{2}f(|g(t_{n+1})|^2)} g(t_{n+1}), \label{known}
\end{eqnarray}
so that the boundary fits perfectly with that to which we arrive advancing $\tau/2$ from $u_n$ with the nonlinear part of Schr\"odinger equation and to that to which we want to arrive when advancing $\tau/2$ from $w^n(\tau)$ with the same nonlinear part.

According to \cite{ACP}, in order to integrate (\ref{midstrang}) without order reduction with a rational-like method which is based on the implicit midpoint rule, for which the stability function is
\begin{eqnarray}
r(z)=-1+\frac{2}{1-\frac{z}{2}},
\label{r}
\end{eqnarray}
the following must be implemented
\begin{eqnarray}
\tilde{w}_{n+1}= r(i \tau A_0) \bar{w}_n- F_n(\tau)g^n(0+\cdot),  \nonumber
\end{eqnarray}
where $\tilde{w}_n=w^n(0)-K g^n(0)$ and
$$F_n(\tau) g^n =(I-i \frac{\tau}{2}A_0)^{-1} K \boldsymbol{\eta}^T  g^n(0+\tau \mathbf{d}),$$
so that
$\boldsymbol{\eta}^T g^n(0+\tau \mathbf{d})$ approximates $L\tau B (I- \tau B/2)^{-1} g^n(0+\cdot)$ with order $3$, where $B$ is the operator consisting of the first time derivative applied over functions in $C^1_{ub}([0,\infty), H^{\frac{3}{2}}(\partial \Omega))$ and $L: C_{ub}([0,\infty),H^{\frac{3}{2}}(\partial \Omega))\to H^{\frac{3}{2}}(\partial \Omega))$ denotes the delta operator
$$L h= h(0), \quad h \in C_{ub}([0,\infty),H^{\frac{3}{2}}(\partial \Omega)).$$
Then, the approximation to $w^n(\tau)$ would be given by $\tilde{w}_{n+1}+ K g^n(\tau)$.

According to Lemma 4.3 in \cite{AP}, $\boldsymbol{\eta}=(\eta_1,\eta_2,\eta_3)^T$ and $\mathbf{d}=(d_1,d_2,d_3)^T$ must be chosen so that
$$
d_1^j\eta_1+d_2^j \eta_2+d_3^j \eta_3=j! F_j, \quad 0 \le j \le 2,
$$
where $F_j$ are the first Taylor coefficients of $F(z)=z/(1-z/2)$ around $z=0$. We notice that, by taking $d_1=0, d_2=1/2, d_3=1$, we do not only manage to consider the known values (\ref{known}) but also that the coefficient $\eta_2=0$ so that the intermediate value $g^n(\tau/2)$ is in fact not necessary. Notice that the system to solve is
\begin{eqnarray}
\begin{array}{ccccccc}
\eta_1&+&\eta_2&+&\eta_3&=&0,  \\
&&  \eta_2/2&+&\eta_3&=&1,  \\
&& \eta_2/4 &+&\eta_3&=&1,
\end{array}
\end{eqnarray}
for which $\eta_1=-1$, $\eta_2=0$, $\eta_3=1$.

Plugging this together,
\begin{eqnarray}
\tilde{w}_{n+1}&=&[-I+2(I-\frac{i \tau}{2}A_0)^{-1}]\tilde{w}_n- (I-\frac{i \tau}{2}A_0)^{-1}K [g^n(\tau)-g^n(0)],
\nonumber
\end{eqnarray}
from what $w_{n+1}=\tilde{w}_{n+1}+K g^n(\tau) \approx w^n(\tau)$ would be
\begin{eqnarray}
%\hspace{1cm}
w_{n+1}&=&-w^n(0)+K [g^n(0)+g^n(\tau)]
\nonumber \\
&&+(I-\frac{i \tau}{2} A_0)^{-1}\bigg[2 w^n(0)- K[g^n(0)+g^n(\tau)]\bigg].
\label{wn1}
\end{eqnarray}
According to \cite{ACP}, the error commited in this way satisfies the following when $w^n \in C^3([0,\tau],L^2(\Omega))$ and $g^n \in C^4([0,\tau],L^2(\partial \Omega))$
\begin{eqnarray}
\hspace{0.5cm}\|w_{n+1}-w^n(\tau)\|_{L^2(\Omega)}\le C \tau^3 (\max_{0\le j \le 3} \|{w^n}^{(j)}(0)\|_{L^2(\Omega)}+\max_{0\le j \le 4} \|{g^n}^{(j)}(0)\|_{L^2(\partial \Omega)}), \label{erlocmp}
\end{eqnarray}
for some constant $C$.



\subsection{Symmetry}
In this subsection, we will prove that advancing $-\tau$ from $w_{n+1}$ using the same formula in (\ref{wn1}), we would arrive at $w^n(0)$.
For that, it suffices to take into account that (\ref{wn1}) can be written as
\begin{eqnarray}
w_{n+1}=[-I +2(I-\frac{i \tau}{2}A_0)^{-1}]w^n(0)+[I-(I-\frac{i \tau}{2}A_0)^{-1}]K [g^n(0)+g^n(\tau)],
\label{sim1}
\end{eqnarray}
and that changing $\tau$ by $-\tau$, $w_{n+1}$ by $w^n(0)$ and $g^n(0)$ by $g^n(\tau)$, it follows that
\begin{eqnarray}
w^n(0)=[-I +2(I+\frac{i \tau}{2}A_0)^{-1}]w_{n+1}+[I-(I+\frac{i \tau}{2}A_0)^{-1}]K [g^n(0)+g^n(\tau)].
\label{sim2}
\end{eqnarray}
Then, by using the algebraic expressions
$$
(-1+\frac{2}{1+z})^{-1}=-1+\frac{2}{1-z}, \quad (-1+\frac{2}{1+z})^{-1}(1-\frac{1}{1+z})=-1+\frac{1}{1-z},$$
(\ref{sim1}) can be seen equivalent to (\ref{sim2}), and the symmetry follows.



%Demostración tras integración en espacio
%In this subsection, we will prove that advancing $-\tau$ from $W_N^{n+1}$ using the same formula in (\ref{formula}), we would arrive at $W_N^n$. For that, we change $n$ by $n+1$ and $\tau$ by $-\tau$ in (\ref{formula}) and check that it holds whenever (\ref{formula}) holds. More precisely, by doing that,
%\begin{eqnarray}
%W_N^n=-W_N^{n+1} +(I+\frac{i \tau}{2} A_{N,0})^{-1}\bigg[2 W_N^{n+1}-\frac{i \tau}{2} C_N[g^n(0)+g^n(\tau)]\bigg],
%\nonumber
%\end{eqnarray}
%from what
%$$(I+\frac{i \tau}{2}A_{N,0}) W_h^n= (I-\frac{i \tau}{2}A_{N,0}) W_N^{n+1}-\frac{i \tau}{2} C_N[g^n(0)+g^n(\tau)],$$
%and, by applying $(I- \frac{i \tau}{2} A_{N,0})^{-1}$ at both sides, and using that
%$$
%I+\frac{i \tau}{2}A_{N,0}=2 I-(I-\frac{i \tau}{2} A_{N,0}),$$
%we again arrive at (\ref{formula}).



\section{Modified Strang method}



In the following, we will denote the modified Strang method to that which consists of the calculation of $w_{n+1}$ through (\ref{wn1}) starting from $w^n(0)$ in (\ref{midstrang}) and then completing the step through the formula
\begin{eqnarray}
u_{n+1}=e^{i \frac{\tau}{2}f(|w_{n+1}|^2)} w_{n+1}. \label{finalu}
\end{eqnarray}
\subsection{Local error}
Third local order can be achieved in such a way, as the following theorem states for $\rho_n=\hat{u}_{n+1}-u(t_{n+1})$, where $\hat{u}_{n+1}$ is that obtained through (\ref{midstrang}),(\ref{wn1}) and (\ref{finalu}) when starting from $u_n=u(t_n)$.

\begin{theorem}
Let us assume that the solution $u$ of (\ref{nlse}) belongs to
\begin{eqnarray}
\quad C^3([0,T], L^2(\Omega))\cap C([0,T], H^6(\Omega))  \cap C^1([0,T], H^4(\Omega)) \cap C^2([0,T],L^{\infty}(\Omega)),
\label{reg2}
\end{eqnarray}
where $f\in C^2(\mathbb{R}, \mathbb{R})$ and $g\in C^2([0,T], H^{\frac{3}{2}}(\partial \Omega))$. Then, for small enough $\tau$,
\begin{eqnarray}
\|\rho_n\|_{L^2(\Omega)} &\le& C \tau^3 (\|u(t_n)\|_{H^6(\Omega)}+\|\dot{u}(t_n)\|_{H^4(\Omega)}
+\max_{0\le j \le 3} \|u^{(j)}(t_n)\|_{L^2(\Omega)}
\nonumber \\
&&\hspace{6cm}+\max_{0\le j \le 2}\|g^{(j)}(t_n)\|_{H^{\frac{3}{2}}(\partial \Omega)}),
\label{cotaloc}
\end{eqnarray}
for some constant $C$, which does not depend on $\tau$ either on $t_n$.
\label{thloc}
\end{theorem}

\begin{proof}
For the sake of simplicity, let us denote
\begin{eqnarray}
H_n=e^{i \frac{\tau}{2}f(|u(t_n)|^2)}u(t_n), \quad H_{n+1}=e^{-i \frac{\tau}{2}f(|u(t_{n+1})|^2)}u(t_{n+1}).
\label{HH}
\end{eqnarray}
Then, we consider $w^{n,b}(s)$ the polynomial of second degree in $s$ which takes the value $H_n$ at $s=0$, $H_{n+1}$ at $s=\tau$ and a value $X_n$ to be determined at $s=\tau/2$. More precisely,
$$
w^{n,b}(s)=H_n+\frac{H_{n+1}-H_n}{\tau}s+\frac{2(H_{n+1}+H_n-2 X_n)}{\tau^2}s(s-\tau).$$
Let us then denote by $\hat{w}^n$ the solution of (\ref{midstrang}) with initial condition $\hat{w}^n(0)=H_n$ and $g^n(s)=\partial w^{n,b}(s)$, which obviously satisfies (\ref{known}). We notice that $\hat{w}^n(s)-w^{n,b}(s)$ satisfies the following
\begin{eqnarray}
\quad \stackrel{\cdot}{\overbrace{(\hat{w}^n-w^{n,b})}}(s)&=&i A (\hat{w}^n(s)-w^{n,b}(s))+i A H_n-\frac{4 X_n-3 H_n-H_{n+1}}{\tau} \nonumber \\
\quad &&+s \big[\frac{i}{\tau}A(-H_{n+1}-3 H_n+4 X_n)-\frac{4}{\tau^2}(H_{n+1}+H_n-2 X_n)\big] \nonumber \\
\quad &&+\frac{2 i s^2}{\tau^2}A(H_{n+1}+H_n-2 X_n), \nonumber \\
\quad (\hat{w}^n-w^{n,b})(0)&=&0, \nonumber \\
\quad \partial (\hat{w}^n-w^{n,b})(s)&=&0. \label{diferencia}
\end{eqnarray}
Then, if we take $X_n$ such that the second term of the right-hand side of the first line in (\ref{diferencia}) vanishes, i.e.
$$X_n=\frac{1}{4}[ i \tau A H_n+H_{n+1}+3 H_n],$$
it can be deduced that
\begin{eqnarray}
H_{n+1}+H_n-2 X_n&=&\frac{H_{n+1}-H_n}{2}-\frac{i \tau}{2} A H_n. \nonumber
\end{eqnarray}
We then notice that, by Taylor expansions in powers of $\tau$ and using (\ref{nlse}),
\begin{eqnarray}
H_n&=&u+\frac{i \tau}{2}f(|u|^2)u-\frac{\tau^2}{8}f(|u|^2)^2 u- \frac{i \tau^3}{48} f(|u|^2)^3 u +O(\tau^4), \nonumber \\
H_{n+1}-H_n&=& i \tau A u+\frac{\tau^2}{2}[\ddot{u}-i f(|u|^2) \dot{u}-i f'(|u|^2)(u \dot{\bar{u}}+\bar{u}\dot{u})u]+\tau^3 D, \nonumber
\end{eqnarray}
where, for the sake of brevity, we have suppressed the evaluation at $t_n$ and where
\begin{eqnarray}
D&=&\frac{1}{6}\stackrel{\dots}{u}-\frac{i}{4}f'(|u|^2)(\bar{u}\ddot{u}+2 \dot{u}\dot{\bar{u}}+u \ddot{\bar{u}})u-\frac{1}{4}f(|u|^2)f'(|u|^2)(u\dot{\bar{u}}+\bar{u}\dot{u})u \nonumber \\
&&-\frac{i}{2}f'(|u|^2)(u \dot{\bar{u}}+\bar{u} \dot{u})\dot{u}-\frac{1}{8}f(|u|^2)^2 \dot{u}-\frac{i}{4}f(|u|^2) \ddot{u} \nonumber \\
&&+\frac{i}{24}f(|u|^2)^3 u -\frac{i}{4} f''(|u|^2)(u \dot{\bar{u}}+\bar{u}\dot{u})^2 u +O(\tau). \nonumber
%\\&=&-\frac{\tau^2}{4}A^2 u(t_n)+O(\tau^3), \nonumber
\end{eqnarray}
Using this, (\ref{diferencia}) can be written as
\begin{eqnarray}
\stackrel{\cdot}{\overbrace{(\hat{w}^n-w^{n,b})}}(s)&=&i A (\hat{w}^n(s)-w^{n,b}(s)) \nonumber \\
&&+s \big[-\frac{3 i \tau}{4}A^2 \big( f(|u(t_n)|^2) u(t_n)\big)-4 \tau D(t_n) +O(\tau^2)\big] \nonumber \\
&&+ s^2 \big[ -\frac{i}{2} A^3 u (t_n)+O(\tau) \big], \nonumber \\
(\hat{w}^n-w^{n,b})(0)&=&0, \nonumber \\
\partial (\hat{w}^n-w^{n,b})(s)&=&0. \label{diferenciab}
\end{eqnarray}
By using then the variation-of-constants formula and the definition of the $\varphi_j$-functions \cite{HO} , which are well-known to imply that
\begin{eqnarray}
\varphi_j( t A_0)=\frac{1}{t^j} \int_0^t
e^{(t-\tau)A_0}\frac{\tau^{j-1}}{(j-1)!}d\tau, \quad j \ge 1,
%\label{varphi}
\nonumber
\end{eqnarray}
it follows that
\begin{eqnarray}
\hat{w}^n(s)-w^{n,b}(s)&=&s^2 \tau \varphi_2(i s A_0)\big[-\frac{3 i}{4}A^2 \big( f(|u(t_n)|^2) u(t_n)\big)-4 D(t_n) +O(\tau)\big]
\nonumber \\
 &&+s^3  \varphi_3(i s A_0)[-\frac{i}{2} A^3 u(t_n)+O(\tau)].
\label{exprphi}
\end{eqnarray}
Therefore, as $\|\varphi_2(i s A_0)\|_{L^2(\Omega)}$ and $\|\varphi_3(i s A_0)\|_{L^2(\Omega)}$ are bounded,
\begin{eqnarray}
\|\hat{w}^n(\tau)-H_{n+1}\|_{L^2(\Omega)}&=&\|\hat{w}^n(\tau)-w^{n,b}(\tau)\|_{L^2(\Omega)} \nonumber \\
&\le& C \tau^3 (\|u(t_n)\|_{H^6(\Omega)}+\|\dot{u}(t_n)\|_{H^4(\Omega)}+\max_{0\le j \le 3} \|u^{(j)}(t_n)\|_{L^2(\Omega)}),
\nonumber
\end{eqnarray}
where we have used that, from (\ref{nlse}),
$$A^2 \big( f(|u|^2)u\big)=-i A^2 \dot{u}-A^3 u,$$
and the assumed regularity (\ref{reg2}).

Considering now $\hat{w}_{n+1}$ as the approximation to $\hat{w}^n(\tau)$ given by the rational-like midpoint rule without order reduction, and using (\ref{erlocmp}) and the fact that
\begin{eqnarray}
&&\hat{w}^n(0)=H_n, \, \hat{w}^{n(1)}(0)=i A H_n, \, \hat{w}^{n(2)}(0)=-A^2 H_n, \, \hat{w}^{n(3)}(0)=-i A^3 H_n,\nonumber \\
&&g^n(0)=\partial H_n, \,  g^{n(1)}(0)=\partial [i A H_n],
\, g^{n(2)}(0)=-\partial A^2 u(t_n)+O(\tau), \, g^{n(j)}(0)=0, \, j=3,4, \nonumber
\end{eqnarray}
it happens that
\begin{eqnarray}
\|\hat{w}_{n+1}-H_{n+1}\|_{L^2(\Omega)}&\le& \|\hat{w}_{n+1}-\hat{w}^n(\tau)\|_{L^2(\Omega)}+\|\hat{w}^n(\tau)-H_{n+1}\|_{L^2(\Omega)}
\nonumber \\
&\le & C \tau^3 (\|u(t_n)\|_{H^6(\Omega)}+\|\dot{u}(t_n)\|_{H^4(\Omega)}+\max_{0\le j \le 3} \|u^{(j)}(t_n)\|_{L^2(\Omega)}
\nonumber \\
&& \hspace{2cm}+ \max_{0\le j \le 2}\| \partial A^j u(t_n)\|_{H^{\frac{3}{2}}(\partial \Omega)}).
\label{cot1}
\end{eqnarray}
We now notice that
\begin{eqnarray}
|\hat{w}_{n+1}|^2-|H_{n+1}|^2=(\hat{w}_{n+1}-H_{n+1}) \bar{\hat{w}}_{n+1}+H_{n+1}(\bar{\hat{w}}_{n+1}-\bar{H}_{n+1}),
\label{for1}
\end{eqnarray}
where both $H_{n+1}$ and $\bar{\hat{w}}_{n+1}$ belong to $L^\infty(\Omega)$ because of (\ref{reg2}),(\ref{acotK2}) and Hille-Yoshida theorem using that, from (\ref{wn1}),
\begin{eqnarray}
\hat{w}_{n+1}&=&-H_n+K [\partial H_n+\partial H_{n+1}]+(I-\frac{i \tau}{2} A_0)^{-1}\bigg[2 H_n- K[\partial H_n+\partial H_{n+1}]\bigg].
\nonumber
\end{eqnarray}
Therefore,
$$
\| |\hat{w}_{n+1}|^2-|H_{n+1}|^2\|_{L^2(\Omega)}\le C \|\hat{w}_{n+1}-H_{n+1}\|_{L^2(\Omega)},$$
which implies that
$$
\| f(|\hat{w}_{n+1}|^2)-f(|H_{n+1}|^2)\|_{L^2(\Omega)}\le C \|\hat{w}_{n+1}-H_{n+1}\|_{L^2(\Omega)},$$
and thus, using that for any $\theta_1$ and $\theta_2$,
\begin{eqnarray}
|e^{i \theta_1}-e^{i \theta_2}|\le |\theta_1-\theta_2|,
\label{theta12}
\end{eqnarray}
it follows that
\begin{eqnarray}
\| e^{i \frac{\tau}{2}f(|\hat{w}_{n+1}|^2)}-e^{i \frac{\tau}{2}f(|H_{n+1}|^2)}\|_{L^2(\Omega)}\le C \tau \|\hat{w}_{n+1}-H_{n+1}\|_{L^2(\Omega)}.
\label{cot2}
\end{eqnarray}
Then, using (\ref{HH}),
\begin{eqnarray}
\hat{u}_{n+1}&=&e^{i \frac{\tau}{2} f(|\hat{w}_{n+1}|^2)} \hat{w}_{n+1}
\nonumber \\
&=& e^{i \frac{\tau}{2} f(|\hat{w}_{n+1}|^2)}(\hat{w}_{n+1}-H_{n+1})+\big[e^{i \frac{\tau}{2}f(|\hat{w}_{n+1}|^2)}-e^{i \frac{\tau}{2}f(|H_{n+1}|^2)} \big]H_{n+1}+u(t_{n+1}),  \nonumber
\end{eqnarray}
from what (\ref{cotaloc}) follows using (\ref{cot1}) and the fact that $\partial A^j u(t_n)$ and $g^{(j)}(t_n)$ are closely related through (\ref{nlse}), and then considering (\ref{cot2}) and that $e^{i \frac{\tau}{2}f(|\hat{w}_{n+1}|^2)}$ and $H_{n+1}$ belong to $L^\infty(\Omega)$.
\end{proof}


\begin{thebibliography}{99}
\bibitem[ACP]{ACP} C. Arranz-Sim\'on, B. Cano \& C. Palencia, {\it Rational methods for abstract linear initial boundary value problems without order reduction}, accepted in Numerical Algorithms. https://doi.org/10.48550/arXiv.2510.15481
\bibitem[AP]{AP} C. Arranz-Sim\'on \& C. Palencia, {\it Rational methods for abstract linear non-homogeneous problems without order reduction}, SIAM J. Num. Anal. \textbf{63}(1) (2025), 422--436.
\bibitem[CR]{CR} B. Cano \& N. Reguera, {\it Avoiding order reduction when integrating nonlinear Schr\"odinger equation with Strang method}, J. Comp. Appl. Math. \textbf{316} (2017), 86--99.
\bibitem[HO]{HO} M. Hochbruck \& A. Ostermann, {\it Exponential Integrators}, Acta Numerica (2010), 209--286.
\bibitem[S]{S} J. C. Strikwerda, {\it Finite difference schemes and partial differential equations}, Pacific Grove California: Wadsworth \& Brooks, 1989.
\bibitem[WH]{WH} J. A. C. Weideman \& B. M. Herbst, {\it Split-step methods for the solution of the nonlinear Schr\"odinger equation}, SIAM J. Numer. Anal. \textbf{23} (1986), 485--507.
\end{thebibliography}
\end{document}
